\hspace{3mm}

\centerline{\bf \huge 6 класс}

\begin{flushright}
    \scriptsize{ $-$ Если сдаваться, не попробовав, результат не изменится.}
\end{flushright}

\problem{1} Расул Владимирович придумал систему перевода цифр в буквы и обратно, в которой каждой из десяти цифр соответствует одна уникальная буква. Расул Владимирович решил поделиться своей системой с директором Элистинского технического лицея. Директор же решила воспользоваться интересной системой и сложила в столбик два слова ЭЛМШ с ЭТЛ, получив суммарно МШМЭЛ. Определите все варианты чисел, скрывающихся за ЭЛМШ с ЭТЛ, и докажите, что других вариантов быть не может.

\hspace{3mm}

\problem{2} Однажды лососи и шиншиллы выстроились в один ряд. Известно, что среди любых двадцати подряд идущих поровну лососей и шиншилл, а среди любых двадцати двух подряд идущих не поровну лососей и шиншилл. Какое максимальное количество животных могло выстроиться в один ряд?

\hspace{3mm}

\problem{3} В турнире по настольному теннису участвовали $512$ человек. В каждом туре участники разбивались на пары и играли, после чего проигравшие выбывали из турнира (ничьих не бывает). После турнира, когда определился единоличный победитель, каждый из участников заявил, что в течение турнира обыграл не более одного правши. Причём все правши сказали правду, а все левши соврали. Сколько правшей могло участвовать в турнире?

\hspace{3mm}

\newpage

\hspace{3mm}


\problem{4} После олимпиады Герман решил подшутить над Расулом Владимировичем и Анджуром Аркадьевичем и подговорил всех учеников ЭлМШ прийти на апелляцию. Шиншиллы по своей неопытности не умели лить воду, поэтому просто молчали на апелляции. Лососи из-за недостаточного количества опыта лили воду любому проверяющему. Киви же, зная учителей, лили воду только Анджуру Аркадьевичу. Оказалось, что олимпиаду проверили просто идеально, проставив всем только заслуженные баллы. Но Анджур Аркадьевич все равно добавлял по баллу тем, кто уверенно лил воду. Расул Владимирович же, наоборот, отнимал балл у тех, кто пытался лить воду. После апелляции оказалось, что баллов стало на 15 больше. Сколько Шиншилл апеллировалось у Расула Владимировича, если количество попавших к нему Киви было столько же, сколько Шиншилл попало к Анджуру Аркадьевичу. Причем Анджур Аркадьевич и Расул Владимирович приняли одинаковое количество детей.

\hspace{3mm}


\problem{5} Божья коровка и муравей устроили забег на часах. Муравей бежал по минутной стрелке часов, а божья коровка - по часовой стрелке. Когда впервые с начала гонки стрелки на часах пересеклись, муравей пробежал $\dfrac{2}{5}$ дистанции, а божья коровка только $\dfrac{1}{4}$. Когда божья коровка завершила забег, если муравей добежал до конца стрелки в 12:00? Забег начался одновременно, а божья коровка и муравей бежали с постоянной скоростью.